% new_document_template.tex
% =============================================================================
%                    TEMPLATE VOOR CURSUSSAMENVATTING
%                    KU Leuven - School Macros v2.0
% =============================================================================
% Dit is een starttemplate voor nieuwe documenten.
% Kopieer dit bestand en pas het aan voor je eigen cursus.
%
% Tip: Compileer TWEEMAAL voor correcte formularium en symbolenlijst!
% =============================================================================

\documentclass[a4paper,11pt]{article}
\usepackage{school-macros}

% =============================================================================
%                           DOCUMENT INFO
% =============================================================================
% Pas deze gegevens aan voor je cursus:
\newcommand{\vaknaam}{Naam van het Vak}
\newcommand{\vakcode}{XXXYYY}
\newcommand{\auteur}{Jouw Naam}
\newcommand{\academiejaar}{2025-2026}

\title{\vaknaam{} --- Samenvatting}
\author{\auteur}
\date{Academiejaar \academiejaar}

% =============================================================================
%                              DOCUMENT
% =============================================================================
\begin{document}

% --- Titelpagina (kies één van beide opties) ---

% Optie 1: Eenvoudige titel
\maketitle

% Optie 2: KU Leuven stijl (uncomment onderstaande, comment \maketitle)
% \kuleuventitle{\vaknaam{} --- Samenvatting}{\vakcode}{\auteur}{Academiejaar \academiejaar}

\tableofcontents
\newpage

% =============================================================================
%                              INHOUD
% =============================================================================

\section{Inleiding}

% Begin hier met je inhoud. Enkele tips:
% - Gebruik \concept{} voor nieuwe begrippen
% - Gebruik \frm{}{}{} voor belangrijke formules  
% - Gebruik \sym{}{}{} bij eerste vermelding van symbolen
% - Gebruik \SI{}{} voor getallen met eenheden

\begin{conceptbox}[title=Voorbeeld Definitie]
Dit is een voorbeeld van een conceptbox. Gebruik deze voor belangrijke definities.
\end{conceptbox}

\section{Theorie}

\begin{theorieblok}[title=Voorbeeld Theorie]
Gebruik de theorieblok voor theoretische uitleg en achtergrond.
\end{theorieblok}

% Voorbeeld formule (wordt automatisch toegevoegd aan formularium)
\frm{Voorbeeld Formule}{E = mc^2}{waarbij $m$ de massa [\si{kg}] en $c$ de lichtsnelheid [\si{m/s}]}

\section{Oefeningen}

\begin{oefenblok}[title=Oefening 1]
\textbf{Gegeven:} $m = \SI{2}{kg}$, $v = \SI{10}{m/s}$

\textbf{Gevraagd:} Bereken de kinetische energie.

\textbf{Oplossing:}
\[ E_k = \frac{1}{2}mv^2 = \frac{1}{2} \times \SI{2}{kg} \times (\SI{10}{m/s})^2 = \SI{100}{J} \]
\end{oefenblok}

\begin{examenbox}
Dit is een compacte examentip. Ideaal voor korte, praktische tips.
\end{examenbox}

% =============================================================================
%                              APPENDICES  
% =============================================================================
\newpage
\appendix

% Automatisch gegenereerde lijsten (compileer tweemaal!)
\printsymbols
\printformularium

\end{document}
